\usepackage{xeCJK}
\usepackage{geometry}
\geometry{left=2cm,right=2cm,top=2.5cm,bottom=2.5cm}
\setlength{\headheight}{12.7pt}

\usepackage{titlesec}
\titleformat{\section}
  {\normalfont\large\bfseries}
  {\thesection}
  {1em}
  {}
\titlespacing*{\section}
  {0pt}
  {3.5ex plus 1ex minus .2ex}
  {2.3ex plus .2ex}

\usepackage{amsmath}
\usepackage{amssymb}
\usepackage{amsthm}
\usepackage{enumitem}
\setlist[enumerate]{itemsep=0pt,topsep=0pt,parsep=0pt,leftmargin=0pt,itemindent=2.5em}
\setlist[itemize]{itemsep=0pt,topsep=0pt,parsep=0pt,leftmargin=0pt,itemindent=2.5em}
\usepackage{fancyhdr}
\pagestyle{fancy}
\fancyhf{}
\usepackage{graphicx}
\usepackage{hyperref}
\usepackage{indentfirst}
\usepackage{lmodern}


\begin{document}
\setlength{\abovedisplayskip}{1pt}
\setlength{\belowdisplayskip}{1pt}

\section{Aleph 数}

\hypertarget{sec:定义1.1}{}
\noindent\textbf{定义1.1 (Aleph数)}

递归定义类映射$\aleph:\mathbf{On}\to\mathbf{Cn}$如下:
\vspace{3pt}
\[
    \forall\alpha\in\mathbf{On}:\quad
    \aleph_\alpha\overset{\mathrm{def}}{=}\left\{\begin{aligned}
        &\omega,\quad&&\alpha=0,\\[-3pt]
        &\aleph(\max\alpha)^{\mathrm{h}},\quad&&\alpha\text{\,\,是后继序数},\\[-3pt]
        &\sup\aleph[\alpha],\quad&&\alpha\text{\,\,是极限序数}.
    \end{aligned}\right.
\]

\vspace{10pt}

\hypertarget{sec:定理1.1}{}
\noindent\textbf{定理1.1 (Aleph数)}

\begin{enumerate}
    \item $\aleph_0=\omega$,
    \item 任给序数$\alpha$, $\aleph_{\alpha+1}=\aleph_\alpha^{\mathrm{h}}$,
    \item 任给序数$\alpha$, 如果$\alpha$是极限序数, 那么
    $\displaystyle\aleph_\alpha=\sup_{\beta<\alpha}\aleph_\beta$.
\end{enumerate}

\vspace{15pt}

显然任给序数$\alpha$, $\aleph_\alpha$都是
无限\href{01 基数#nameddest=sec:定义1.1}{基数}.

\vspace{10pt}

\hypertarget{sec:定理1.2}{}
\noindent\textbf{定理1.2}

任给序数$\alpha,\alpha'$, 如果$\alpha<\alpha'$, 那么$\aleph_\alpha<\aleph_{\alpha'}$.

\noindent{\kaishu 证明.}

给定$\alpha$, 对$\alpha'>\alpha$归纳证明$\aleph_{\alpha'}>\aleph_\alpha$.
任取$\alpha'>\alpha$, 假设
\[
    \forall\gamma:\quad\alpha<\gamma<\alpha'\to\aleph_\gamma>\aleph_\alpha.
\]

\begin{enumerate}
    \item 如果$\alpha'$是后继序数, 即有$\alpha'=\gamma+1$,
    那么有$\alpha\leqslant\gamma<\alpha'$, 依假设有
    \vspace{1pt}
    \[
        \aleph_{\alpha'}=\aleph_\gamma^{\mathrm{h}}>
        \aleph_\gamma\geqslant\aleph_\alpha.\\[3pt]
    \]
    \item 如果$\alpha'$是极限序数, 则$\alpha+1<\alpha'$, 从而
    \vspace{1pt}
    \begin{equation*}
        \aleph_{\alpha'}=\sup_{\beta<\alpha'}\aleph_\beta
        \geqslant\aleph_{\alpha+1}=\aleph_\alpha^{\mathrm{h}}>\aleph_\alpha.
    \tag*{\qedsymbol}\end{equation*}
\end{enumerate}

\vspace{10pt}

\hypertarget{sec:定理1.3}{}
\noindent\textbf{定理1.3}

任给序数$\alpha$, 有$\alpha\leqslant\aleph_\alpha$.

\noindent{\kaishu 证明.}

对$\alpha$归纳证明$\alpha\leqslant\aleph_\alpha$. 任取$\alpha$,
假设对任意的$\beta<\alpha$有$\beta\leqslant\aleph_\beta$.

\begin{enumerate}
    \item 如果$\alpha=0$, 显然$0\leqslant\omega=\aleph_0$.
    \item 如果$\alpha$是后继序数, 即有$\alpha=\beta+1$,
    那么$\beta<\alpha$, 依假设有
    \[
        \alpha=\beta+1\leqslant\aleph_\beta+1\leqslant
        \aleph_\beta^{\mathrm{h}}=\aleph_\alpha.
    \]
    \item 如果$\alpha$是极限序数, 那么依假设有
    \begin{equation*}
        \alpha=\sup_{\beta<\alpha}\beta\leqslant
        \sup_{\beta<\alpha}\aleph_\beta=\aleph_\alpha.
    \tag*{\qedsymbol}\end{equation*}
\end{enumerate}





\newpage

最后, Aleph数可以表示所有的无限基数.

\vspace{10pt}

\hypertarget{sec:定理1.4}{}
\noindent\textbf{定理1.4}

任给无限基数$\kappa$, 存在唯一的序数$\alpha$使得$\kappa=\aleph_\alpha$.

\noindent{\kaishu 证明.}

\noindent{\kaishu 存在性.}

取满足$\aleph_\theta>\kappa$的最小的$\theta$,
即对任意的$\alpha<\theta$有$\aleph_\alpha\leqslant\kappa$.
由$\kappa$无限得$\theta$非零.

\vspace{3pt}
如果$\theta$是极限序数, 那么$\displaystyle
\aleph_\theta=\sup_{\alpha<\theta}\aleph_\alpha\leqslant\kappa$矛盾.
所以$\theta$是后继序数, 有$\theta=\alpha+1$.

\vspace{3pt}
于是由$\alpha<\theta$得$\aleph_\alpha\leqslant\kappa$.
假设$\kappa>\aleph_\alpha$则有
\[
    \kappa\geqslant\aleph_\alpha^{\mathrm{h}}=\aleph_{\alpha+1}=\aleph_\theta,
\]
矛盾. 所以$\kappa=\aleph_\alpha$.

\noindent{\kaishu 唯一性.} 因为$\aleph$保序, 显然. \hfill $\qedsymbol$

\vspace{30pt}

\hypertarget{sec:定义1.2}{}
\noindent\textbf{定义1.2 (后继基数与极限基数)}

任给无限基数$\kappa$, 存在唯一的序数$\alpha$使得$\kappa=\aleph_\alpha$ :
\begin{enumerate}
    \item 如果$\alpha$是后继序数, 那么称$\kappa$是一个\textbf{后继基数};
    \item 如果$\alpha$是$0$或极限序数, 那么称$\kappa$是一个\textbf{极限基数}.
\end{enumerate}

\vspace{10pt}

另外, 无限基数都是极限序数.

\vspace{10pt}

\hypertarget{sec:定理1.5}{}
\noindent\textbf{定理1.5}

无限基数都是极限序数.

\noindent{\kaishu 证明.}

假设无限基数$\kappa$是后继序数, 即有$\kappa=\alpha+1$. 显然$\alpha\geqslant\omega$,
令$\varepsilon=\alpha-\omega$. 于是
\[
    \kappa=\alpha+1=\omega+\varepsilon+1\approx 1+\omega+\varepsilon=
    \omega+\varepsilon=\alpha,
\]
但是由$\alpha<\kappa$得$\alpha\prec\kappa$, 矛盾. \hfill $\qedsymbol$

\end{document}